\section{Fermat型の固定Poissonスキームと相対Poissonモジュライの既約成分}
\label{sec:pg-fermat-relative-components}
本節では、\emph{固定した滑らかなcubicのPoisson scheme}と、
\emph{cubic自身も変形させる相対Poisson scheme}を比較する。
いずれの「既約成分」もスキームの基礎位相空間（すなわち被約化）の既約成分を意味し、
その非被約構造や成分交叉のイデアルについての主張ではない。
Poisson双ベクトルのスカラーは割らず、相対側でのみ射影座標変換を割る。
本節の分類は定理\ref{thm:pg-normal}に依存し、その証明全体の独立再監査・
先行研究上の新規性確認は完了していない。

\subsection{対象と二つの成分数}
$N=n+2\ge5$、$V=\mathbb C^N$、$S=\operatorname{Sym}(V^*)$とする。
$F\in S_3$が滑らかな射影cubic $X_F\subset\mathbb P(V)$を定めるとき、
命題\ref{prop:pg-represent}により、複素点の集合として
\[
 \operatorname{Pois}(X_F)(\mathbb C)=
 \{P\in\Lambda^3V:\,[\iota_{dF}P,\iota_{dF}P]=0\}
\]
と同一視する（右辺には射影化の規約を適用する）。
Jacobi方程式は$P$について二次斉次である。
$F_N=\sum_{i=1}^N z_i^3$、$X_N=V(F_N)$について
$\mathfrak P_N^{\mathrm{Fer}}=\operatorname{Pois}(X_N)$と書く。

また滑らかな三次式の射影パラメータ開集合を
$U_N\subset\mathbb P(S_3)$とし、その上の相対Poissonスキームの非零部分を
$\mathfrak I_N^\times$とする。$[F]$の局所表示を選ぶと、
双ベクトルは$\pi=\iota_{dF}P$と書ける。
$F$のスカラー変更には$P$を逆スカラー倍して対応するので、
相対の定義は局所表示の選択に依存しない。
相対モジュライは商スタック
$\mathfrak M_n=[\mathfrak I_N^\times/\mathrm{PGL}(V)]$の基礎位相空間として扱う。

$\left\{\begin{smallmatrix}N\\k\end{smallmatrix}\right\}_{\ge3}$を、
$\{1,\ldots,N\}$を各部分集合の要素数が$3$以上である
$k$個の\emph{非順序}非空部分集合へ分割する方法の数とする。

\begin{theorem}[Fermat fiberと相対schemeの成分数]
\label{thm:pg-fermat-relative-count}
$n\ge3$、$N=n+2$とする。定理\ref{thm:pg-normal}の下で、
\begin{enumerate}
\item Fermat cubic $X_N$のアフィンPoissonスキームの被約既約成分は
\[
 \#\operatorname{Irr}(\mathfrak P_N^{\mathrm{Fer}})
 =\sum_{k=1}^{\lfloor N/3\rfloor}
  \left\{\begin{smallmatrix}N\\k\end{smallmatrix}\right\}_{\ge3}.
 \tag{F}\label{eq:pg-fermat-count}
\]
 一般Poisson階数$2k$の成分の個数は
 $\left\{\begin{smallmatrix}N\\k\end{smallmatrix}\right\}_{\ge3}$で、
 各成分のアフィン次元は$3N-8k$である。
\item cubic自身も動かす相対スキームでは、一般階数$2k$ごとに
 ただ一つの既約成分が存在し、従って
\[
 \#\operatorname{Irr}(\mathfrak M_n)=\left\lfloor\frac{N}{3}\right\rfloor.
 \tag{R}\label{eq:pg-relative-count}
\]
 その商スタック次元は$m=N-3k$として
\begin{equation}
 \label{eq:pg-relative-dim}
 d_{n,k}=k(3N-8)+\binom{m+2}{3}
 +k\left(10+6m+3\binom{m+1}{2}\right)-N^2.
\end{equation}
\end{enumerate}
$\pi=0$を追加しても双方の被約既約成分数は変わらない。
\end{theorem}

\subsection{固定Fermat fiberの成分を座標分割で記述する}
\begin{proof}[定理\ref{thm:pg-fermat-relative-count}の第一項]
まず$P\in\Lambda^3V$が$F_N$に関してPoissonであるとする。
FermatのHessianは
$\operatorname{Hess}_{F_N}(z)=6\operatorname{diag}(z_1,\ldots,z_N)$である。
補題\ref{lem:pg-linearize}のJacobi条件を対角rank一二次形式に適用すると
\begin{equation}
 \label{eq:pg-fermat-contractions}
 (\iota_{dz_i}P)\wedge(\iota_{dz_i}P)=0
 \quad(1\le i\le N)
\end{equation}
を得る（これは各縮約二形式のPl\"ucker方程式そのものである）。

定理\ref{thm:pg-normal}により、$P=\sum_{s=1}^k P_s$、
$P_s\in\Lambda^3 W_s\setminus\{0\}$、$\dim W_s=3$、
$W_s$は互いに独立と書ける。
各$\iota_{dz_i}P_s$は$\Lambda^2W_s$に属する二形式である。
もし同じ$i$に対して二つの異なる$s,t$の縮約が非零ならば、
それらの外積は$\Lambda^2W_s\otimes\Lambda^2W_t$の非零元であり、
他の組の寄与と相殺できないから\eqref{eq:pg-fermat-contractions}に反する。
従って、各座標線形形式$dz_i$は高々一つの$W_s$上でしか非零にならない。
$B_s=\{i:dz_i|_{W_s}\ne0\}$とすれば、$B_s$は互いに素で、
$|B_s|\ge3$かつ$W_s\subset V_{B_s}:=\langle e_i:i\in B_s\rangle$である。
未使用の座標添字をいずれかの$B_s$へ加えれば、
$\{1,\ldots,N\}$の大きさ$3$以上の$k$個のブロックへの分割$\mathcal B$を得る。

逆に、そのような分割$\mathcal B=\{B_1,\ldots,B_k\}$に対して
\[
 Z_{\mathcal B}=\bigoplus_{s=1}^k
 \widehat{\operatorname{Gr}}(3,V_{B_s})
 \ \subset\ \Lambda^3V
\]
を考える。ハットはPl\"ucker埋込のアフィン錐（零点も含む）を表す。
これは線形独立な$\Lambda^3V_{B_s}$の中の錐の直積であり、閉かつ既約である。
各$P_s$は分解可能で、$F_N=\sum_s\sum_{i\in B_s}z_i^3$は
それぞれ別の座標に依存するから、各Jacobian Poisson括弧はJacobi恒等式を満たし、
異なるブロック間の括弧も零である。従って
$Z_{\mathcal B}\subset\operatorname{Pois}(X_N)$である。
先の議論からこれらの$Z_{\mathcal B}$は被約集合全体を被覆する。

$|B_s|=b_s$とすると
$\dim\widehat{\operatorname{Gr}}(3,V_{B_s})=3(b_s-3)+1$なので
\[
 \dim Z_{\mathcal B}=\sum_s(3b_s-8)=3N-8k.
\]
異なる$k$については、縮約写像
$\Lambda^2V^*\to V$のrankが一般点で$3k$なので、
小さい$k$の成分の閉包は大きい$k$の一般点を含めない。
逆向きの包含も上の次元が$k$とともに厳密に減少するので不可能である。
同じ$k\ge2$の異なる分割については、各$W_s$がそのブロックの
全座標方向に非零射影を持つ一般点を取る。
このとき縮約二形式の平方零集合から$W_s$の組が順序を除いて
内在的に復元される（補題\ref{ne:blocks}）。
したがって各$W_s$の座標支持$B_s$も復元され、別の分割による
$Z_{\mathcal B'}$へその一般点は入らない。
$k=1$の分割はただ一つである。
ゆえに全$Z_{\mathcal B}$が相互に異なる既約成分で、
式\eqref{eq:pg-fermat-count}が得られる。
零点は全ての錐に含まれ、追加の成分にはならない。
\end{proof}

\subsection{動くcubicのrank層は一つずつになる}
\begin{proof}[定理\ref{thm:pg-fermat-relative-count}の第二項]
$P_k=e_{123}+e_{456}+\cdots+e_{3k-2,3k-1,3k}$と置く。
定理\ref{thm:pg-normal}によって、rank$\,2k$の相対Poisson対の三重ベクトル部分は
$P_k$の$\mathrm{GL}_N$軌道$\mathcal O_{N,k}$である。
この軌道は既約で、その次元は$k(3N-8)$である。
実際$N=3k+m$とすると、$P_k$の安定化子のLie次元は
$8k+3km+m^2$なので、$N^2-(8k+3km+m^2)=k(3N-8)$となる。
$P_k$を固定した三次式は
\[
 F=c(u)+\sum_{s=1}^k\bigl(f_{s,3}(x_s)+f_{s,2}(x_s,u)+f_{s,1}(x_s,u)\bigr)
\]
の形の線形空間であり、係数次元は
\[
 \ell_{N,k}=\binom{m+2}{3}
 +k\left(10+6m+3\binom{m+1}{2}\right).
\]
その射影化の滑らかなcubicの部分は非空既約開集合である
（Fermat三次式が例になる）。よってrank$\,2k$層は既約で
次元$k(3N-8)+\ell_{N,k}-1$を持ち、
$\dim\mathrm{PGL}_N=N^2-1$を引くと式\eqref{eq:pg-relative-dim}を得る。

異なるrank層の閉包が互いに包含されないことを確かめる。
小さい$k$の層の閉包では縮約写像のrankが高々$3k$という閉条件が残る。
一方、隣り合う層について直接計算すると
\[
 \bigl(k(3N-8)+\ell_{N,k}-1\bigr)
 -\bigl((k+1)(3N-8)+\ell_{N,k+1}-1\bigr)
 =3(3k-1)(N-3k)+8>0.
\]
従って大きい$k$の層の閉包が小さい$k$の層全体を含むこともない。
以上によりrankごとに一つずつの既約成分となる。
零切断はrank二の相対錐の閉包に含まれる。
\end{proof}

\subsection{数値の比較と意味}
分割数の指数母関数は
\begin{equation}
\label{eq:pg-fermat-egf}
 \sum_{N\ge0}\frac{t^N}{N!}
 \left\{\begin{smallmatrix}N\\k\end{smallmatrix}\right\}_{\ge3}
 =\frac1{k!}(e^t-1-t-t^2/2)^k.
\end{equation}
従って係数抽出だけで固定Fermatの各rank成分数を求められる。
\begin{center}
\begin{tabular}{r|rrrr|r|r}
\toprule
$n$ & Fermat rank 2 & rank 4 & rank 6 & rank 8
& Fermat計 & 相対計 \\
\midrule
3 & 1 & -- & -- & -- & 1 & 1 \\
4 & 1 & 10 & -- & -- & 11 & 2 \\
5 & 1 & 35 & -- & -- & 36 & 2 \\
6 & 1 & 91 & -- & -- & 92 & 2 \\
7 & 1 & 210 & 280 & -- & 491 & 3 \\
8 & 1 & 456 & 2100 & -- & 2557 & 3 \\
9 & 1 & 957 & 10395 & -- & 11353 & 3 \\
10 & 1 & 1969 & 42735 & 15400 & 60105 & 4 \\
\bottomrule
\end{tabular}
\end{center}
例えば$n=7$（$N=9$）では、rank四の$210$は
$(3,6)$型の分割$\binom93=84$と$(4,5)$型の分割$\binom94=126$の和、
rank六の$280$は$(3,3,3)$型の分割
$9!/((3!)^3\,3!)$であり、$1+210+280=491$となる。
一方、相対的にはこの三つのrankにつき一つずつの三成分だけが残る。
これは\emph{固定特殊ファイバーの既約成分が、全体の異なる既約成分と
一対一に対応する必要がない}ことを具体的に示している。

\begin{remark}[最大性・scheme構造・既知結果との区別]
固定Fermatの成分数が、全ての滑らかなcubicの中で最大になることは
上の議論からは従わない。特に高次元での最大性は別の問題である。
表は\emph{被約既約成分数}についてのもので、Jacobi方程式で定義される
Poisson schemeが被約であるとか、成分の交叉や零点での局所環を
完全に決定したものではない。
また$\mathfrak M_n$は商スタックである。
その粗モジュライの存在・構成や全ての幾何学的性質をここで主張しない。
式\eqref{eq:pg-fermat-count}--\eqref{eq:pg-relative-dim}の初出・新規性は
包括的には未確認であり、正規形の既知入力と本節の計数的帰結を区別する。
\end{remark}
